% ---- pg-498 （书页 486） ----

\[
\begin{aligned}
\mathbf{W}&=\sum_{i=1}^{3}\sum_{j=1}^{n_i}(\mathbf{x}_{ij}-\bar{\mathbf{x}}_i)(\mathbf{x}_{ij}-\bar{\mathbf{x}}_i)'=(n_1+n_2+n_3-3)\mathbf{S}_{\mathrm{p}}\\
&=\begin{bmatrix}6&-2\\-2&24\end{bmatrix}
\end{aligned}
\]
\[
\mathbf{W}^{-1}=\frac{1}{140}\begin{bmatrix}24&2\\2&6\end{bmatrix};\qquad
\mathbf{W}^{-1}\mathbf{B}=\begin{bmatrix}0.3571&0.4667\\0.0714&0.9000\end{bmatrix}
\]
为求 $\mathbf{W}^{-1}\mathbf{B}$ 的 $s\le\min(g-1,p)=\min(2,2)=2$ 个非零特征值, 必须解方程
\[
\left|\mathbf{W}^{-1}\mathbf{B}-\lambda\mathbf{I}\right|=\left|\begin{bmatrix}0.3571-\lambda&0.4667\\0.0714&0.9000-\lambda\end{bmatrix}\right|=0
\]
或
\[
(0.3571-\lambda)(0.9000-\lambda)-0.4667\times 0.0714=\lambda^2-1.2571\lambda+0.2881=0
\]
用二次求根公式, 得到 $\hat{\lambda}_1=0.9556,\hat{\lambda}_2=0.3015$. 求解方程
\[
(\mathbf{W}^{-1}\mathbf{B}-\hat{\lambda}_i\mathbf{I})\hat{\mathbf{a}}_i=\mathbf{0}\qquad i=1,2
\]
并取解的长度满足 $\hat{\mathbf{a}}_i'\mathbf{S}_{\mathrm{p}}\hat{\mathbf{a}}_i=1$, 可以得到两个标准化特征向量 $\hat{\mathbf{a}}_1,\hat{\mathbf{a}}_2$. 例如方程
\[
(\mathbf{W}^{-1}\mathbf{B}-\hat{\lambda}_1\mathbf{I})\hat{\mathbf{a}}_1=\begin{bmatrix}0.3571-0.9556&0.4667\\0.0714&0.9000-0.9556\end{bmatrix}\begin{bmatrix}\hat{a}_{11}\\\hat{a}_{12}\end{bmatrix}=\begin{bmatrix}0\\0\end{bmatrix}
\]
的解经过标准化 $\hat{\mathbf{a}}_1'\mathbf{S}_{\mathrm{p}}\hat{\mathbf{a}}_1=1$ 之后为 $\hat{\mathbf{a}}_1'=[0.386,0.495]$. 类似的, $\hat{\mathbf{a}}_2'=[0.938,-0.112]$. 两个判别量为
\[
\hat{y}_1=\hat{\mathbf{a}}_1'\mathbf{x}=[0.386\quad 0.495]\begin{bmatrix}x_1\\x_2\end{bmatrix}=0.386x_1+0.495x_2
\]
\[
\hat{y}_2=\hat{\mathbf{a}}_2'\mathbf{x}=[0.938\quad -0.112]\begin{bmatrix}x_1\\x_2\end{bmatrix}=0.938x_1-0.112x_2
\]
\hfill$\blacksquare$

\noindent\textbf{例 11.14}\quad (原油数据的费希尔判别量)

格利尔德和兰茨(Gerrild \& Lantz, 见文献[13])从加利福尼亚埃尔克山的砂岩中采集了一组原油样本. 这些原油可以分入以下三类之一:

$\pi_1$: 威赫尔姆砂岩

$\pi_2$: 下马里尼亚砂岩

$\pi_3$: 高层砂岩

为说明问题, 我们仅考虑五个变量:
\[
\begin{aligned}
X_1&=\text{钒(灰分中的百分比)}\\
X_2&=\sqrt{\text{铁(灰分中的百分比)}}\\
X_3&=\sqrt{\text{铍(灰分中的百分比)}}\\
X_4&=1/\text{饱和烃(面积百分比)}\\
X_5&=\text{芳香烃(面积百分比)}
\end{aligned}
\]
前三个变量是痕量元素, 后两个变量则根据气体色谱化学分析法所产生的一段曲线来确定. \hyperref[tab:11.7]{表 11.7}(见练习 11.30)给出五个原始变量(钒、铁、铍、饱和烃与芳香烃)在 56 个样本中的值. 这些样本的分类是已知的.
